% Sección 1: el problema y qué es KERALYTIX, en lenguaje simple

\seccion{El proyecto}

% ===============================================================
% 1. El problema
% ===============================================================
\begin{frame}{El problema: detectar el queratocono a tiempo}
    % Beamer apoya el bloque de columnas sobre el margen izquierdo (no lo centra).
    % El \hspace* corre el texto hasta el borde de la caja de \idea.
    % totalwidth tiene que alcanzar para el corrimiento + las dos columnas:
    % 0,032 + 0,48 + 0,44 = 0,952 <= 0,96
    \begin{columns}[T, totalwidth=0.96\textwidth]
        \hspace*{0.032\textwidth}
        \begin{column}{0.48\textwidth}
            \vspace{0.25 cm}
            \textbf{Características de la patología}
            \begin{itemize}
                \item La córnea se adelgaza y se deforma hasta tomar forma de cono.
                \item Aparece en jóvenes y avanza rápido. Al principio se confunde con miopía o astigmatismo.
                \item Incluso con IA, 1 de cada 10 casos iniciales pasa desapercibido.
                \item Detectada tarde, puede terminar en un trasplante de córnea.
            \end{itemize}
        \end{column}
        \begin{column}{0.44\textwidth}
            % Una debajo de la otra. Cada ilustración ya trae su rótulo
            % ("Córnea sana" y "Córnea con queratocono")
            \centering
            \includegraphics[height=2.8cm]{queratocono_deformacion_a.png}\par\vspace{0.1cm}
            \includegraphics[height=2.8cm]{queratocono_deformacion_b.png}
        \end{column}
    \end{columns}
    \idea{El problema no es la enfermedad, es detectarla a tiempo.}
\end{frame}

\begin{frame}{Qué mide el equipo}
    \centering
    \resizebox{\linewidth}{!}{% Esquema de los parámetros corneales (Marco teórico)
% (a) corte de una córnea con queratocono sobre sus esferas de referencia
% (b) vista de frente: ubicación del punto más delgado y orientación en grados
\begin{tikzpicture}[
    font=\small,
    >={Stealth[length=1.8mm]},
    referencia/.style = {dashed, gray!70, thick},
    superficie/.style = {very thick, black!80},
    cota/.style = {<->, thick},
    nota/.style = {font=\footnotesize, align=left, inner sep=1pt,
                   execute at begin node={\setstretch{1}}},
    guia/.style = {thin},
    declare function={
        esfAnt(\x) = sqrt(36 - \x*\x) - 3;          % esfera de referencia anterior
        esfPos(\x) = sqrt(27.04 - \x*\x) - 2.8;     % esfera de referencia posterior
        cono(\x)   = exp(-((\x - 0.8)^2) / 1.6);    % protrusión del cono
        ant(\x)    = esfAnt(\x) + 0.6*cono(\x) - 0.1;
        pos(\x)    = esfPos(\x) + 0.9*cono(\x) - 0.1;
    }
  ]

  % ---------------- (a) Corte de la córnea ----------------
  % Tejido corneal entre ambas superficies
  \fill[blue!10]
    plot[domain=-4:4, samples=120] (\x, {ant(\x)})
    -- plot[domain=4:-4, samples=120] (\x, {pos(\x)}) -- cycle;
  % Esferas de referencia (forma ideal)
  \draw[referencia] plot[domain=-4.2:4.2, samples=80] (\x, {esfAnt(\x)});
  \draw[referencia] plot[domain=-4.2:4.2, samples=80] (\x, {esfPos(\x)});
  % Superficies reales
  \draw[superficie] plot[domain=-4:4, samples=120] (\x, {ant(\x)});
  \draw[superficie] plot[domain=-4:4, samples=120] (\x, {pos(\x)});
  \node[nota] at (-3.2, 2.75) {cara anterior};
  \node[nota] at (-2.0, 1.2) {cara posterior};

  % Zona de 3 mm centrada en el punto más delgado (A y B)
  \draw[orange!80!black, thick, |-|] (-0.15, 4.15) -- (1.75, 4.15)
    node[midway, above=2pt, nota, text=orange!80!black] {zona de 3 mm (A y B)};
  \draw[orange!80!black, densely dotted] (-0.15, 4.15) -- (-0.15, {ant(-0.15)});
  \draw[orange!80!black, densely dotted] (1.75, 4.15) -- (1.75, {ant(1.75)});

  % K max: punto de mayor curvatura de la cara anterior
  \fill[violet] (0.8, {ant(0.8)}) circle (2.2pt);
  \node[nota, text=violet] (lkmax) at (-3.1, 4.3) {\textbf{K max}: punto de mayor\\curvatura de la cara anterior};
  \draw[guia, violet] (lkmax.east) -- (0.75, {ant(0.8)+0.05});

  % Espesor mínimo (C)
  \draw[cota, teal!70!black] (0.8, {pos(0.8)}) -- (0.8, {ant(0.8)});
  \node[nota, text=teal!70!black] (lesp) at (5.6, 3.7) {\textbf{Espesor mínimo (C)}:\\espesor en el punto\\más delgado};
  \draw[guia, teal!70!black] (lesp.west) -- (0.85, {(ant(0.8)+pos(0.8))/2});

  % Elevaciones respecto de las esferas de referencia
  \draw[cota, red!70!black] (1.75, {esfAnt(1.75)}) -- (1.75, {ant(1.75)});
  \draw[cota, red!70!black] (1.75, {esfPos(1.75)}) -- (1.75, {pos(1.75)});
  \node[nota, text=red!70!black] (lelev) at (5.6, 2.1) {\textbf{Elevación} anterior\\y posterior: cuánto\\sobresale cada cara\\de su esfera};
  \draw[guia, red!70!black] (lelev.west) -- (1.8, {(esfAnt(1.75)+ant(1.75))/2});
  \draw[guia, red!70!black] (lelev.west) -- (1.8, {(esfPos(1.75)+pos(1.75))/2});

  % Leyenda
  \draw[referencia] (-4, 0.15) -- ++(0.6, 0) node[right, nota, text=black] {esfera de referencia (forma ideal)};
  \fill[blue!10] (-4, -0.45) rectangle ++(0.6, 0.3);
  \node[nota, right] at (-3.3, -0.3) {tejido corneal};
  \node[font=\small\bfseries] at (0.6, -1.05) {(a) Corte de una córnea con queratocono};
  \node[nota, text=gray] at (0.6, -1.5) {Esquema sin escala: las deformaciones están exageradas.};

  % ---------------- (b) Vista de frente ----------------
  \begin{scope}[shift={(10.6, 2.0)}]
    \fill[blue!6] (0,0) circle (2.1);
    \draw[thick] (0,0) circle (2.1);
    \draw[gray!60] (-2.35,0) -- (2.35,0);
    \draw[gray!60] (0,-2.35) -- (0,2.35);
    \node[nota] at (2.65, 0) {0°};
    \node[nota] at (-2.8, 0) {180°};
    \node[nota] at (0, 2.6) {90°};
    % Dos orientaciones casi iguales: 1° y 179°
    \draw[red!70!black, thick] (0,0) -- (1:2.1) node[pos=0.72, above=1pt, nota, text=red!70!black] {1°};
    \draw[red!70!black, thick, dashed] (0,0) -- (179:2.1) node[pos=0.72, above=1pt, nota, text=red!70!black] {179°};
    % Punto más delgado, su ubicación (x, y) y la zona de 3 mm
    \coordinate (pmd) at (0.55, -0.85);
    \draw[orange!80!black, thick] (pmd) circle (0.7);
    \draw[teal!70!black, densely dotted] (pmd) -- (0.55, 0) node[above, nota, text=teal!70!black] {$x$};
    \draw[teal!70!black, densely dotted] (pmd) -- (0, -0.85) node[left, nota, text=teal!70!black] {$y$};
    \fill[teal!70!black] (pmd) circle (2.2pt);
    \node[nota, align=center] at (0, -2.85) {\textcolor{teal!70!black}{punto más delgado en $(x, y)$}\\\textcolor{orange!80!black}{y zona de 3 mm a su alrededor}};
    \node[nota, align=center, text=red!70!black] at (0, 3.15) {1° y 179° marcan casi\\la misma dirección};
    \node[font=\small\bfseries] at (0, -3.05) [below=8pt] {(b) Vista de frente};
  \end{scope}
\end{tikzpicture}
}
\end{frame}

% ===============================================================
% 2. Qué es KERALYTIX (diagrama de bloques general)
% ===============================================================
\begin{frame}{Qué es KERALYTIX}
    Una herramienta web que le da al oftalmólogo una segunda opinión sobre el estudio.
    \vspace{0.1cm}
    \begin{center}
        \begin{tikzpicture}[node distance=0.45cm and 0.9cm]
            \node[caja=3.2cm] (mapas) {\textbf{Mapas de la córnea}\\ \footnotesize los cuatro que exporta el equipo de topografía};
            \node[caja=3.2cm, below=of mapas] (clin) {\textbf{Datos del paciente}\\ \footnotesize edad, sexo, curvatura, espesor};
            \coordinate (medio) at ($(mapas.east)!0.5!(clin.east)$);
            \node[caja=3.2cm, fill=azul!15, very thick, right=1.2cm of medio, anchor=west] (modelo)
                {\textbf{Modelo de IA}\\ \footnotesize analiza imágenes y datos juntos};
            \node[caja=3.2cm, draw=verde, right=of modelo] (informe)
                {\textbf{Informe}\\ \footnotesize normal o queratocono, confianza y mapa de calor};
            \draw[flecha] (mapas.east) -- ++(0.6cm,0) |- (modelo.west);
            \draw[flecha] (clin.east) -- ++(0.6cm,0) |- (modelo.west);
            \draw[flecha] (modelo) -- (informe);
            \node[font=\footnotesize\bfseries, text=gray, above=0.15cm of mapas] {Qué recibe};
            \node[font=\footnotesize\bfseries, text=gray] at (modelo |- {$(mapas.north)+(0,0.4cm)$}) {Qué hace};
            \node[font=\footnotesize\bfseries, text=gray] at (informe |- {$(mapas.north)+(0,0.4cm)$}) {Qué entrega};
        \end{tikzpicture}
    \end{center}
    \footnotesize{Los mapas son imágenes y los datos del paciente son información tabulada. Al combinar ambos tipos, el modelo es \textbf{multimodal}. El mapa de calor marca qué zonas de la córnea llevaron al resultado. La decisión final es siempre del médico.}
\end{frame}

% ===============================================================
% 3. Objetivos
% ===============================================================
\begin{frame}{Objetivos}
    \begin{block}{Objetivo general}
        \small Desarrollar un sistema de soporte a la decisión clínica, basado en IA sobre mapas corneales y datos clínicos,
        para la detección temprana del queratocono y accesible mediante una aplicación web.
    \end{block}
    \vspace{0.3cm}
    \begin{tabla}{colspec = {Q[c,m,0.8cm] X[l,m]}, rows = {font=\small}}
        N.º & Objetivos específicos \\
        1 & Analizar un conjunto de datos público con mapas corneales y variables clínicas. \\
        2 & Construir un modelo de IA que combine mapas y variables para identificar el queratocono. \\
        3 & Integrar una explicación visual de las zonas de la córnea que llevaron a cada resultado. \\
        4 & Implementar una aplicación web para cargar estudios y consultar el informe. \\
    \end{tabla}
\end{frame}

% ===============================================================
% 4. Fases y estado actual
% ===============================================================
\begin{frame}{Fases del proyecto y estado actual}
    \begin{center}
        \begin{tikzpicture}[x=1.9cm, y=1cm]
            % Meses
            \foreach \m/\n in {1/Ago, 2/Sep, 3/Oct, 4/Nov, 5/Dic}
                \node[font=\footnotesize, text=gray] at (\m-0.5, 3.3) {\n};
            % Fases
            \fill[verde!70] (0,2.2) rectangle (1,2.8);
            \node[anchor=east, font=\small, align=right] at (-0.1,2.5) {Fase 1: comprensión y\\selección de datos};
            \fill[azul!70] (1,1.2) rectangle (4,1.8);
            \node[anchor=east, font=\small] at (-0.1,1.5) {Fase 2: modelado};
            \fill[gray!40] (3,0.2) rectangle (5,0.8);
            \node[anchor=east, font=\small] at (-0.1,0.5) {Fase 3: aplicación web};
            % Objetivos que cubre cada fase
            \node[font=\scriptsize, text=white] at (0.5,2.5) {Obj. 1};
            \node[font=\scriptsize, text=white] at (2.5,1.5) {Obj. 2 y 3};
            \node[font=\scriptsize] at (4,0.5) {Obj. 4};
            % Hoy
            \draw[rojo, very thick, dashed] (3.3,0) -- (3.3,3.05);
            \node[rojo, font=\footnotesize\bfseries, anchor=south] at (3.3,-0.55) {hoy};
        \end{tikzpicture}
    \end{center}
    \vspace{0.1cm}
    \centering
    {\small\textcolor{verde}{\textbf{$\blacksquare$}} Terminada \quad
     \textcolor{azul}{\textbf{$\blacksquare$}} En curso \quad
     \textcolor{gray!60}{\textbf{$\blacksquare$}} Por empezar}
\end{frame}
